% Chapter 1 draft: Sections 1.1 and 1.2
% Suggested bibliography command in main thesis file:
% \bibliographystyle{apalike}
% \bibliography{references_chapter1}

\section{Alzheimer's Disease: Clinical Relevance and Biological Definition}

Alzheimer's disease (AD) is the most common cause of dementia and one of the major clinical challenges associated with population ageing. It is a progressive neurodegenerative disorder in which cognitive and functional decline usually emerge after a long period of underlying biological change. Clinically, AD is commonly associated with impairment in episodic memory, followed by decline in additional cognitive domains and progressive loss of independence. However, the current understanding of AD is no longer limited to the stage in which dementia is clinically apparent. The disease is increasingly interpreted as a long biological process that precedes, and eventually may lead to, the clinical syndrome of dementia \citep{scheltens2021alzheimer,jack2018nia}.

From a pathological perspective, AD is characterized by the accumulation of amyloid-\(\beta\) plaques and tau neurofibrillary tangles, together with synaptic dysfunction, neuronal injury, and neurodegeneration. These processes do not occur as a single event, nor do they necessarily coincide with the first observable cognitive symptoms. Instead, AD develops across a continuum in which molecular and cellular abnormalities may be detectable before dementia-level impairment becomes evident \citep{scheltens2021alzheimer,sperling2011preclinical}. This temporal separation between biological disease and clinical manifestation is central to modern AD research, because it makes early detection and disease staging conceptually different from the traditional diagnosis of dementia.

A major conceptual shift in the field has been the transition from a purely clinical definition of AD to a biological definition. The 2018 NIA-AA Research Framework explicitly defines AD by its underlying pathologic processes rather than by the presence of clinical symptoms alone \citep{jack2018nia}. In this framework, cognitive status and biological disease status are separated: an individual may show biomarker evidence of AD pathology while remaining cognitively unimpaired, or may present cognitive impairment whose underlying cause requires biological confirmation. The framework organizes AD biomarkers according to the AT(N) system, where ``A'' refers to amyloid-\(\beta\) deposition, ``T'' to pathologic tau, and ``N'' to neurodegeneration or neuronal injury \citep{jack2018nia}.

The revised Alzheimer's Association criteria further reinforce the biological framing of AD and reflect the increasing role of biomarkers in diagnosis and staging \citep{jack2024revised}. This shift is especially important because it clarifies that AD should not be equated automatically with dementia. Dementia is a clinical syndrome, whereas AD is a biological disease process that can be present across different clinical stages. At the same time, many public biomedical datasets still use diagnostic labels such as AD, mild cognitive impairment (MCI), and cognitively normal control, often without complete biomarker confirmation. Therefore, in computational studies, these labels should be interpreted as practical supervised learning targets, not as perfect representations of underlying biological disease states.

This distinction is directly relevant to the present thesis. The classification task considered here focuses on distinguishing AD from MCI using blood gene expression data. This is not equivalent to separating diseased individuals from healthy controls. Instead, it involves diagnostic categories that may lie close to one another along the clinical and biological continuum of AD. For this reason, the task requires a cautious interpretation: a model trained on diagnostic labels may capture differences related to disease stage, symptom severity, systemic biological alterations, or dataset-specific effects, rather than a direct and isolated signature of AD pathology.

\section{Mild Cognitive Impairment and the Alzheimer's Disease Continuum}

Mild Cognitive Impairment is a clinical diagnostic construct introduced to describe cognitive impairment that lies between normal ageing and dementia. Petersen defines MCI as an early but abnormal state of cognitive impairment, distinct from normal cognitive ageing and not severe enough to meet criteria for dementia \citep{petersen2004mci}. In this view, MCI occupies an intermediate zone along the cognitive continuum, where the distinction from normal ageing can be subtle and the transition toward very early dementia can also be clinically challenging \citep{petersen2004mci}. This makes MCI a useful but intrinsically heterogeneous category.

The clinical characterization of MCI depends on the presence of measurable cognitive impairment with relatively preserved functional independence. In the amnestic form, the classical criteria include a memory complaint, preferably corroborated by an informant, objective memory impairment for age, essentially preserved general cognitive function, largely intact activities of daily living, and absence of dementia \citep{petersen2004mci}. These criteria are clinically important because they identify individuals whose cognitive performance is abnormal but whose impairment does not yet reach the threshold required for dementia. However, Petersen also emphasizes that the diagnosis relies partly on clinical judgement, because there is no single universally accepted cutoff separating normal ageing, MCI, and very early dementia \citep{petersen2004mci}.

MCI should not be treated as synonymous with AD. It includes different clinical subtypes, including amnestic MCI, multiple-domain MCI, and non-amnestic single-domain impairment, and these subtypes may reflect different underlying etiologies \citep{petersen2004mci}. Amnestic MCI, especially when presumed to have a degenerative cause, is often considered a possible prodromal form of AD. Other forms of MCI, however, may progress toward non-AD dementias or may reflect vascular, psychiatric, traumatic, or mixed causes \citep{petersen2004mci}. This heterogeneity is essential for the interpretation of AD versus MCI classification, because the MCI group may contain individuals at different biological stages and with different disease mechanisms.

The NIA-AA workgroup recommendations for MCI due to AD refine this distinction by introducing the concept of MCI as a symptomatic predementia phase that may be attributable to AD when supported by clinical course and biomarker evidence \citep{albert2011diagnosis}. In this framework, the diagnosis of MCI due to AD is not simply based on cognitive impairment, but on the likelihood that the impairment reflects underlying AD pathophysiology. This likelihood can be strengthened by biomarkers related to amyloid deposition, neuronal injury, or neurodegeneration \citep{albert2011diagnosis}. Therefore, MCI due to AD is more specific than the broader clinical category of MCI, but this level of specificity is not always available in public transcriptomic datasets.

The concept of preclinical AD further extends the continuum model. Sperling and colleagues argue that the pathophysiological process of AD may begin years before the diagnosis of dementia and before clear clinical impairment is present \citep{sperling2011preclinical}. This implies that the clinical categories of cognitively normal, MCI, and AD dementia do not map perfectly onto biological disease status. Some cognitively normal individuals may already have AD pathology, some MCI individuals may be on an AD trajectory, and some MCI individuals may have impairment unrelated to AD. The 2018 NIA-AA biological framework makes this separation explicit by treating biomarker-defined disease status and clinical impairment as related but distinct dimensions \citep{jack2018nia}.

For this thesis, the AD versus MCI task is therefore best understood as a supervised approximation of a clinically and biologically complex boundary. It separates individuals diagnosed with AD from individuals diagnosed with MCI, but it does not directly separate biological AD from non-AD unless biomarker confirmation is available. This is a stronger and more realistic framing than treating the problem as a simple binary disease classification task. It also motivates the use of representation-learning approaches, because blood gene expression may reflect subtle systemic differences associated with disease stage, neurodegeneration, inflammation, or other biological processes, while also being affected by the high-dimensional and heterogeneous nature of transcriptomic data.
